% Cover letter for IEEE Signal Processing Letters submission
% Compile: pdflatex cover_letter
\documentclass[11pt]{letter}
\usepackage[T1]{fontenc}
\usepackage{lmodern}
\usepackage[margin=2cm]{geometry}
\usepackage{parskip}
\usepackage{enumitem}
\usepackage{hyperref}
\hypersetup{hidelinks}
\setlength{\parskip}{4pt}

\signature{Do Phuc Hao}
\address{Do Phuc Hao\\ Faculty of Information Technology\\ Da Nang Architecture University\\
Da Nang, Vietnam\\ haodp@dau.edu.vn}

\begin{document}

\begin{letter}{Editor-in-Chief\\ IEEE Signal Processing Letters}

\opening{Dear Editor-in-Chief,}

I am pleased to submit the manuscript titled ``Geometric Equivariance as an
Inductive Bias for Orientation-Robust Direction Finding in Non-Terrestrial
Networks'' for consideration as a Letter in IEEE Signal Processing Letters.

The work studies how a known geometric symmetry of a sensor array, the rotation of
the array, should enter a learned direction-of-arrival estimator. On a uniform
circular array, azimuth rotation reduces exactly to a cyclic shift of the sensor
index, the cyclic group $C_M$, which makes the symmetry explicit and lets us compare
three options: ignore it, augment for it with random cyclic shifts, or build it into
a $C_M$-equivariant circular-convolution network. The main findings are concise and,
to my knowledge, new for radio-frequency array processing:

\begin{itemize}[leftmargin=1.4em,itemsep=2pt,topsep=2pt]
\item Built-in equivariance and cyclic-shift augmentation both reach the classical
accuracy on the clean single-source task, but equivariance is more data efficient,
more stable across seeds, and lighter. This controlled equivariance-versus-augmentation
comparison for direction finding is not reported in a literature that obtains
robustness by training over perturbations.
\item On a two-source single-snapshot task the learned models resolve sources about
twice as close as the classical Bartlett beamformer, and the equivariant model alone
resolves the hardest separation, with an order of magnitude fewer parameters than the
augmented baseline. This is where learning exceeds, rather than inherits, classical
behavior.
\end{itemize}

The manuscript fits the scope and concise format of Signal Processing Letters: a
clear, reproducible signal-processing result motivated by non-terrestrial networks,
where antenna arrays sit on platforms whose orientation varies. Prior deep-learning
direction-of-arrival methods use linear arrays and obtain robustness by data
augmentation rather than by an architectural symmetry constraint, and the only
genuinely equivariant direction-finding networks are acoustic, so the cyclic-group
equivariant formulation on a circular array, the equivariance-versus-augmentation
comparison, and the orientation-generalization study are not pre-empted by existing
work.

I confirm that this manuscript is original, has not been published before, and is not
under consideration for publication elsewhere. There is a single author and there are
no conflicts of interest to declare. All numbers and figures are produced by
self-contained Python scripts that are released for full reproducibility.

Thank you for your time and consideration. I look forward to the reviewers' feedback.

\closing{Sincerely,}

\end{letter}
\end{document}